\documentclass[11pt,a4paper]{article}
\usepackage{fontspec}
\setmainfont{Noto Serif}
\setsansfont{Noto Sans}
\usepackage{babel}
\babelprovide[main,import]{russian}
\usepackage{amsmath,amssymb,mathtools,booktabs,geometry,fancyhdr,microtype,xcolor,hyperref}
\geometry{left=21mm,right=21mm,top=21mm,bottom=21mm}
\definecolor{ink}{HTML}{173D57}
\hypersetup{colorlinks=true,urlcolor=ink,linkcolor=ink,pdftitle={BGS-F. Совпадающие участки квантовых траекторий}}
\pagestyle{fancy}\fancyhf{}\lhead{\small BGS-F}\rhead{\small Совпадающие участки квантовых траекторий}\cfoot{\thepage}
\setlength{\headheight}{15pt}\setlength{\parindent}{1em}\setlength{\parskip}{2.5pt}
\setlength{\emergencystretch}{3em}\tolerance=3000
\newcommand{\Tr}{\operatorname{Tr}}
\newcommand{\EE}{\mathbb E}
\newcommand{\R}{\mathbb R}
\newcommand{\Z}{\mathbb Z}
\newcommand{\C}{\mathbb C}
\newcommand{\si}{\mathrm{Sine}_2}
\newcommand{\eps}{\varepsilon}
\newcommand{\sinc}{\operatorname{sinc}}
\newcommand{\sect}[1]{\section*{\sffamily\color{ink}#1}}
\newcommand{\subt}[1]{\par\medskip\noindent\textbf{#1.}\quad}
\newcommand{\T}{\mathbb T}
\newcommand{\dist}{\operatorname{dist}}
\newcommand{\Ad}{\operatorname{Ad}}
\begin{document}
\noindent{\small\sffamily\color{ink} МАТЕМАТИЧЕСКАЯ ФИЗИКА\hfill 20 НОЯБРЯ 2026}
\begin{center}{\LARGE\sffamily Интерференция на гейзенберговском масштабе}\\[4pt]
{\large Кубический и квадратный корни в сумме по парам путей}
\end{center}
Для исходной карты $U=e^{i10^{-80}N A}M$ рассмотрим два замкнутых слова $x,y$ длины $m$ по алфавиту $\Z/N\Z$. Их вклад в $|\Tr U^m|^2$ равен $a(x)\overline{a(y)}$, где $a(x)=\prod_t U_{x_{t+1},x_t}$. Временные индексы берутся по модулю $m$. Слово имеет отмеченное начало, поэтому все $N^m$ слов входят в след.

\sect{Результаты}
При $m/N\to t>0$ разрешим менять одну дугу длиной $L=\lfloor m^{1/3}\rfloor$ с любым началом и совмещать слова любым циклическим сдвигом. Внутри этого класса пары с разными мультимножествами направленных рёбер имеют суммарный вклад $-mt^2/6+o(m)$; пары с одинаковыми рёбрами дают $m(1+t^2/6)+o(m)$. Каждая пара учитывается один раз. При $L\sim b\sqrt N$ обе суммы достигают порядка $N^{3/2}$ с явными коэффициентами. Для длины $D$ кратчайшего участка, допускающего перестановку, доказан закон $\Pr(D>u(N^2/m)^{1/3})\to e^{-u^3/6}$ при выборе пути с весом $|a(x)|^2/\Tr B^m$, $B_{ij}=|U_{ij}|^2$. Доказательства даны в разделах 6а--6з.

Зафиксируем $d>2/3$, $d\le1$. Включим пару один раз, если при некотором $s$ совпадают не менее $dm$ позиций $x_t=y_{s+t}$ или $x_t=y_{s-t}$. Для каждого фиксированного $0<a<b$ существуют $c,C>0$ такие, что по всем достаточно большим нечётным $N$ равномерно при $aN\le m\le bN$
\[
 \boxed{\displaystyle
 \left|\sum_{\rm включённые\ пары}a(x)\overline{a(y)}-m\right|
          \le C e^{-c m}.}                                                   \tag{1}
\]
Если из этой суммы удалить все точные циклические копии, её модуль не более $C'e^{-c'm}$. Для семейства, определённого \emph{только обратными} сравнениями, вся сумма по однократно учтённым парам меньше $e^{-L m}$ по модулю при любом заранее заданном $L>0$ и достаточно большом $N$. Время здесь растёт линейно по $N$. Отдельные слагаемые комплексны.

При одном заранее выбранном сдвиге (1) имеет более сильный аналог для любой доли $d>0$: главный член равен 1 для прямого сравнения и 0 для обратного. Перенос этого последнего порога на объединение событий потребовал бы контроля многократных совпадений. Порог $2/3$ в (1) обеспечивается отдельной оценкой пересечений.

Получены и два вывода о собственных векторах. Пусть $u_j$ --- произвольный ортонормированный собственный базис $U$, $J_{xj}=1/N$. Тогда в обычной операторной норме матриц размера $N$ по всем нечётным $N$
\[
 \big\|( |u_j(x)|^2)_{xj}-J\big\|\longrightarrow0,\qquad
 \big\|( u_j(x)^2)_{xj}\big\|\longrightarrow0.                              \tag{2}
\]
В частности, даже позиционные фазы, выбранные с учётом полного спектра, не могут придать положительной доле собственных состояний почти вещественный вид в одном общем базисе. Точная формулировка через среднее перекрытие дана в разделе 7.

\textbf{Статус.} Для (P) остаётся суммарный вклад пар с меньшим числом совпадений при всех сдвигах в обоих направлениях. Его предел и закон $\mathrm{Sine}_2$ не установлены.
\newpage
\sect{1. Два точных оператора и счёт совпадений}
Введём $V_+=U\otimes\bar U$, $V_-=U\otimes U^*$, проектор $P$ на $\mathrm{span}\{|x,x\rangle\}$ и $Q=I-P$. Разница между покомпонентным сопряжением и сопряжённо-транспонированной матрицей существенна. Раскрытие следа даёт полином
\[
 Z_m^\pm(z)=\Tr\{[(Q+zP)V_\pm]^m\}=\sum_{h=0}^m c_{m,h}^\pm z^h .
\]
Коэффициент --- сумма амплитуд пар с ровно $h$ совпадениями при данном сдвиге. Для знака минус в обратном слове ставим $z_t'=y_{s-t}$, тогда во втором множителе на шаге стоит $\overline{U_{z'_t,z'_{t+1}}}$. Циклическая перенумерация делает коэффициенты независимыми от $s$. При $z=1$ оба полинома равны $|\Tr U^m|^2$.

Положим $e=N^{-1/2}\sum_x |x,x\rangle$. Имеем $V_+e=e$, $Pe=e$; в $Z_m^+$ эта прямая даёт ровно $z^m$. На её ортогональном дополнении используем $V=V_+|_{e^\perp}$, $P_0=P-ee^*$. В обратном случае оставляем $V=V_-$, $P_0=P$. Ниже $K_{j,N}=P_0 V^jP_0$ рассматривается на $\mathrm{ran}\,P_0$. В прямом случае это сжатие на константно-ортогональные векторы матрицы $|(U^j)_{xy}|^2$; при продолжении нулём на константы оно равно $|(U^j)_{xy}|^2-J$. В обратном случае
\[
                    K_{j,N}(x,y)=(U^j)_{xy}\overline{(U^j)_{yx}} .
\]

\textbf{Лемма о числе совпадений.} Пусть $R>1$ и $(R+1)\sum_{j\ge1}\limsup_N\|K_{j,N}\|<1$. Для каждого $r>1$ при больших $N$ спектральный радиус $[I+(z-1)P_0]V$ не превосходит $r$ равномерно по $|z|=R$. Действительно, при $|\lambda|\ge r$ невозмущённый резольвентный ряд сходится. Возможное собственное значение потребовало бы необратимости
\[
 I-(z-1)\sum_{j\ge1}\lambda^{-j}K_{j,N}.
\]
Это следует из формулы конечного ранга либо непосредственно из уравнения на собственный вектор после применения $P_0 V(\lambda-V)^{-1}$. При фиксированном $r$ хвост ограничен $\sum_{j>J}r^{-j}$, поскольку $V$ унитарен. Сначала выбираем $J$, затем $N$; норма поправки строго меньше единицы.

След степени ограничен размерностью, умноженной на $r^m$, в том числе для ненормальной матрицы. Неравенство Коши для коэффициентов при $h\ge\lceil dm\rceil$ даёт
\[
 \sum_{h\ge dm}|c_{m,h}^\pm-\mathbf1_{\pm=+,\,h=m}|
       \le \frac{N^2}{1-R^{-1}}\,(r/R^d)^m .                                \tag{3}
\]
У прямого оператора ниже допустимо $R=100$, у обратного --- любое фиксированное $R$. В обоих случаях можно взять $r=R^{d/2}$. Константа порога по $N$ при этом зависит от $d,R$.
\newpage
\sect{2. Отсутствие обращения, сохраняющего координату}
Пусть $F=\phi^\tau\circ B$, $B=\left(\begin{smallmatrix}3&-4\\-2&3\end{smallmatrix}\right)$, $\phi$ --- поток $Y=-J_0\nabla g/(2\pi)$ для физического символа $g$ исходной карты, $J_0=\left(\begin{smallmatrix}0&1\\-1&0\end{smallmatrix}\right)$. Ни при каком $j\ge1$ и вещественной периодической $v(q)$ непрерывная инволюция $R_v(q,p)=(q,v(q)-p)$ не обращает $F^j$, то есть не выполняется $R_v F^j R_v=F^{-j}$.

Из части о жёсткости берём построенное там сопряжение $FH=HB$, $H=\mathrm{id}+h$ и точную оценку
\[
 \|h-\tau h_1\|_\infty\le10^8\tau^2,\qquad
 h_1=(I-B(\cdot\circ B^{-1}))^{-1}Y .
\]
Если бы $R_v$ обращала $F^j$, то $H^{-1}R_vH$ обращала бы $B^j$ и имела гомологическое действие $R_0(q,p)=(q,-p)$. Умножая на $R_0$, получаем коммутирующее с $B^j$ отображение с тождественным гомологическим действием. Оно является сдвигом: периодическая часть поднятия после вычитания константы удовлетворяет $u(B^jx)=B^ju(x)$, откуда исчезает по ограниченности в обе стороны. Следовательно $R_v H=H(t_b R_0)$. Равенство первых координат, проинтегрированное по тору, даёт $b_q=0\bmod1$, затем
\[
                         h_q(q,p)=h_q(q,b_p-p).
\]
Здесь $h_q$ означает первую компоненту векторного поля $h$.

Проверим несовместимость по двум орбитам Фурье. Положим $\Lambda=3+2\sqrt2$, $e_1=(1,0)$, $e_2=(0,1)$. У их $B^\mathsf{T}$-орбит единственные точки в носителе $\widehat g$ --- сами $e_i$ (у других точек вторая координата по модулю не меньше 3). Ряды устойчивой и неустойчивой частей для $h_1$ приведены в разделе 5 части о жёсткости. Левые собственные ковекторы равны $(1/2,\pm1/\sqrt2)$, соответствующие правые имеют первую координату 1. На частотах $(3,-4)=B^\mathsf{T}e_1$ и $(3,4)=B^{-\mathsf{T}}e_1$ первая компонента $h_1$ имеет один и тот же коэффициент $i\widehat g(e_1)/(\sqrt2\Lambda)$ модуля больше $1/200$. На $e_2$ её коэффициент $z=-i\widehat g(e_2)/2=(-4+3i)/2000$. Эти числа используют физический знак $(-1)^{ab}$ в символе.

При $u=e^{2\pi i b_p}$ отражение первой пары частот требует $|u^4-1|\le2\cdot10^8\tau/(1/200)$, а второй пары $|u-\bar z/z|\le2\cdot10^8\tau/|z|$. Между тем $\rho=\bar z/z=(7+24i)/25$ и $|\rho^4-1|^2=451584/390625>1$. Неравенство $|u^4-\rho^4|\le4|u-\rho|$ даёт противоречие при точном $\tau=10^{-80}$. Теорема касается обращений с указанным действием на координату; для подавления обратных обходов требуется именно эта геометрия.
\newpage
\sect{3. Ветви классического отображения и их действия}
Фиксируем $j\ge1$. Пусть $B^j=\left(\begin{smallmatrix}a&-b\\-c&a\end{smallmatrix}\right)$, $b>0$. При замене $p\mapsto-p$ все четыре элемента $DF$ положительны (оценка отклонения от постоянной матрицы меньше $10^{-70}$ в части о давлении). Поэтому $\partial_p(F^j)_q<0$ повсюду. На плоском поднятии по входной координате $y$ и выходной $x$ входной импульс $p=p(y,x)$ определяется однозначно и аналитически, причём $p(y,x+b)=p(y,x)-1$. Выходную координату можно поднять в $x+n$; на торе получаются $b$ ветвей $n\bmod b$. Обозначим их производящие функции $S_n(x,y)=S(x+n,y)$, где $-S_y$ есть входной импульс, а $S_x$ --- выходной. Такие функции существуют на плоскости по симплектичности.

\textbf{Прямое сравнение.} Для $n\ne l\bmod b$ фаза $S_n-S_l$ не является суммой функций одного аргумента даже на открытой площадке. Иначе по аналитичности смешанная производная исчезла бы всюду. Тогда $p(y,x+n-l)-p(y,x)$ не зависит от $x$. Итерация сдвига до кратного $b$ даёт константу $-(n-l)/b$ независимо от $y$. Значит, сопряжение исходного вертикального сдвига на эту ненулевую рациональную величину посредством $F^j$ сохраняет выходную $q$-координату на торе. Сохраняющее площадь отображение, фиксирующее $q$, есть $(q,p)\mapsto(q,p+s(q))$. Конечность его порядка вынуждает $s$ быть константой на окружности. Получили перевод нетривиального сдвига в сдвиг вопреки теореме (14) части о жёсткости.

\textbf{Обратное сравнение.} Ни при каких $n,l$ разность $S_n(x,y)-S_l(y,x)$ не разделяется. При разделении она имела бы вид $f(x)+g_1(y)$ на всей плоскости. Обе производные $f',g_1'$ однопериодичны. В самом деле, у $S$ разность градиентов с градиентом квадратичной производящей функции $B^j$ периодична относительно решётки $(a,1),(-b,0)$ в координатах $(x,y)$ по ковариантности поднятия при целых сдвигах. Перестановка координат сохраняет эту решётку, так как $a^2\equiv1\bmod b$. Квадратичная функция симметрична, а её сдвиги отличаются линейными членами. Следовательно, градиент разности инвариантен, в частности при $(1,a)$ и $(a,1)$. В силу разделения это даёт нужные единичные периоды.

Дифференцирование разности по $x$ даёт сумму выходного импульса прямой ветви и входного импульса переставленной ветви, равную $f'(x)$. По $y$ получается отрицательная сумма двух остальных импульсов, равная $g_1'(y)$. Тем самым на торе $F^j R_{f'} F^j=R_{-g_1'}$: для каждого входа равенство получается выбором соответствующего поднятия $x$, где тождество верно глобально. Обе правые и средние зеркальные карты инволютивны. Обращение равенства даёт $F^{-j}R_{f'}F^{-j}=R_{-g_1'}$, следовательно $R_{f'}$ обращает $F^{2j}$. Это исключено разделом 2. Аргумент использует знак ориентации второго обхода; одно отсутствие трансляционной симметрии здесь недостаточно.
\newpage
\sect{4. Переход от квантовых матриц к классическому ядру}
При фиксированном $j$ обозначим через $L_j$ оператор, переводящий плотность $f(y)$ в выходную плотность координаты $x$ после $F^j$, когда входной импульс равномерен и независим. Его гладкое ядро $k_j(x,y)=\sum_{n\bmod b}|\partial_xp(y,x+n)|$ дважды стохастично. После отождествления сетки с кусочно-постоянными функциями (ортонормированные индикаторы ячеек) имеем сходимость в операторной норме
\[
  (|(U^j)_{XY}|^2)_{XY}\ \longrightarrow L_j,\qquad
  ((U^j)_{XY}\overline{(U^j)_{YX}})_{XY}\ \longrightarrow0 .               \tag{4}
\]
Матрицы продолжаются нулём на ортогональное дополнение ячеек в $L^2(\T)$. Здесь и далее $X,Y$ --- целые индексы, $x=X/N$, $y=Y/N$.

Дадим детали полуклассического перехода, поскольку покомпонентная слабая сходимость не обеспечила бы (4). На координатной площадке ядро $U^j$ имеет вид
\[
 N^{-1/2}\sum_{n\bmod b}c_{n,N} a_n(x,y)e^{2\pi iN S_n(x,y)}
                   +O_j(N^{-3/2}),\qquad
 |c_{n,N}|=1,\quad |a_n|^2=|\partial_xp(y,x+n)|.                          \tag{5}
\]
Возможные аффинные члены при смене поднятия учитываются в фазе. Для нормировки при нечётном $N$ начнём с $S^0(x,y)=-(3x^2-2xy+3y^2)/8$. Точная сумма четырёх ветвей $(4N)^{-1/2}\sum_{n=0}^3e^{2\pi iN S^0(x+n,y)}$ равна $e^{-2\pi i3N/8}M_{XY}$. После выделения непрерывной квадратичной фазы это конечное тождество в восьмых корнях единицы: для $D=3X^2-2XY+3Y^2$ сумма $\sum_{n=0}^3e_8(-2n(3X-Y)-3Nn^2)=2e_8(-3N+ND)$. В архиве оно проверено точно во всех 256 классах остатков.

Дальнейшее построение (5) использует обычный метод стационарной фазы при фиксированном числе шагов [1,2]. Уточним его применимость. В физическом базисе кик до экспоненцирования есть конечная сумма сдвигов $u(x+a/N)$ с множителями $c_{ab}(-1)^{ab}e^{2\pi ib(x+a/(2N))}$. Подстановка $u=N^{-1/2}A(x,y)e^{2\pi iN S(x,y)}$ в уравнение $\dot u=iN A_N u$ даёт $2\pi S_t=g(x,S_x)$, а следующий порядок переносит полуплотность $A$ по характеристикам. В частности, плотность $|A|^2dx$ переносится вдоль ветви. Для проверки транспортного члена при $v=-g_p/(2\pi)$ получается $A_t+v A_x+\frac12(\partial_xv) A=0$, где $\partial_x$ берётся вдоль графика $p=S_x$. Малая длительность кика сохраняет регулярную проекцию: непосредственно после линейной части касательная к образу вертикали имеет ограниченный наклон (в координатах с $p\mapsto-p$ положительный конус, а отношение компонент после $B$ между $2/3$ и $3/4$). Дифференциал промежуточного кика близок к единице, поэтому первая компонента не обращается в нуль.
\newpage
\sect{4а. Стационарная фаза на сетке и оценка нормы}
При следующем применении $M$ используем ту же точную сумму четырёх квадратичных ветвей. С гладким разбиением по промежуточной координате суммирование Пуассона переводит дискретную сумму в интегралы с добавлением целочисленных линейных фаз. Стационарные точки в точности сшивают выходной и входной импульсы по модулю 1. Они невырождены: наклон после кика отрицателен в физических координатах, а вторая производная квадратичной фазы по промежуточной координате равна $-3/4$. Неограниченно большие частоты Пуассона суммируемы после повторного интегрирования по частям. Каждому заданному исходному импульсу соответствует единственная история промежуточных точек на торе. Формула стационарной фазы умножает амплитуду на обратный квадратный корень модуля гессиана; вместе с множителем $1/2$ это даёт по правилу цепочки перенос плотности входного импульса в выходную координату. Фазовый множитель Маслова локально постоянен. Это доказывает главный член (5) индукцией.

Остаток можно вести до любого фиксированного порядка: уравнения для следующих амплитуд линейны с гладкими источниками по разложениям конечных сдвигов, а стационарная фаза допускает тот же порядок. При построении выбирают запас как минимум в две степени $N$ сверх нужного порядка. На сетке переход из равномерной оценки невязки столбца в его $\ell^2$-норму стоит не больше $\sqrt N$, последующие точные пропагаторы унитарны. Интегрирование невязки вдоль кика и конечное число применений $M$ сохраняют запас. Обрезка готового ряда после главного члена даёт указанную равномерную ошибку. Локальные срезки можно вести на поднятиях и суммировать разбиение единицы на промежуточных графиках. Все константы на этом шаге могут зависеть от $j$.

Рассмотрим теперь смешанные слагаемые произведения двух ядер. Одно такое слагаемое есть матрица $T_{XY}=N^{-1}a(x,y)e^{2\pi iN\Phi(x,y)}$. При неразделимой аналитической $\Phi$ её норма стремится к нулю. Действительно, $\|T\|^4\le\Tr[(T^*T)^2]$. Справа возникает четырёхмерная сеточная сумма с нормировкой $N^{-4}$ и фазой
\[
 \Phi(x,y)-\Phi(x,y')-\Phi(x',y)+\Phi(x',y').
\]
Эта фаза не аффинна: её смешанная производная по $x,y$ не тождественно равна нулю. Поэтому множество точек, в которых все первые производные целые, имеет нулевую меру по аналитичности. Целые частоты на компакте ограничены. Вне окрестности сколь угодно малого объёма этого множества гладким разбиением выбираем координату, производная по которой отделена от целых. Одно дискретное суммирование по частям даёт $O(1/N)$ для нормированной суммы в этом направлении: знаменатель $e^{2\pi iN(\Phi_4(t+1/N)-\Phi_4(t))}-1$ отделён от нуля и меняется за шаг на $O(1/N)$. В окрестности используем верхнюю гладкую срезку малого интеграла. Границы площадок обрабатываются разбиением или отбрасыванием узких полос с малой нормой Гильберта--Шмидта. При перемножении локализуем оба аргумента на прямоугольниках, где фаза имеет аналитическое поднятие.

По разделу 3 при обращённом порядке исчезают в норме все слагаемые, при прямом --- все, кроме одинаковых ветвей. Вклад погрешности (5) после умножения равен $O_j(N^{-2})$ на элемент. Оставшаяся матрица $N^{-1}k_j(x,y)$ сходится к $L_j$ даже в норме Гильберта--Шмидта по непрерывности ядра. Это доказывает (4).
\newpage
\sect{5. Суммируемая оценка и сравнение при одном сдвиге}
Для применения (3) важно суммировать нормы $L_j-J$ до бесконечности \emph{после} предельного перехода в (4); здесь $J$ есть проекция на константы. На частотах входной и выходной координат используются диагональные метки Вейля $(0,b)$. В части о жёсткости доказаны оценки (3),(4) для поднятой эволюции: при $n(a,b)=|a|+\sqrt2|b|$ её взвешенная сумма модулей столбца из $v\ne0$ не превосходит $e^{-.8n(v)}$, а на каждом следующем шаге возникает множитель $q=e^{-.8}$. Вес определяется разностью $|a-\sqrt2 b|-|a+\sqrt2 b|$ и равен единице после экспоненцирования на обеих координатных осях. Нулевая мода изолирована. Следовательно,
\[
 \sum_{b'\ne0}|\widehat{L_j}(b',b)|
                   \le e^{-.8\sqrt2|b|}q^{j-1} \quad(b\ne0).                \tag{6}
\]
Для аккуратного перехода фиксируем $b,b',j$: ограничение поднятых путей расширяющимися шарами радиусов $8^i\sqrt N$, $0\le i\le j$, стоит $o(1)$ по экспоненциальной оценке дальнего перехода той же части. Внутри шаров алиасинг при периодизации также стоит $o(1)$. Формула (4) этого отчёта отождествляет предел элемента в дискретных характерах с коэффициентом $L_j$. Сначала суммируем по произвольному конечному набору выходов, затем снимаем ограничение. Так не требуется равномерная по $j$ формула (5).

Для $1\le j\le20$, $1\le|b|\le30$ есть более сильная оценка для $\ell^2$-нормы спроецированного столбца. Чистая линейная карта отправляет $(0,b)$ строго вне оси. Сравнение с чистой эволюцией по телескопической сумме стоит не более $80\tau\,2|b|\,6^{21}$. Действительно, на чистой метке $v$ изменение при одном кике по норме не более $80\tau n(v)$ из формулы коммутатора, а $n(S^iv)\le6^in(v)$. Для фиксированных меток и $j$ берём $N\to\infty$. Норма оператора ограничена суммой норм столбцов. Для $|b|>30$ и для $j>20$ используем (6). Так как $q<1/2$, $e^{-.8\sqrt2}<1/3$, получаем
\[
 \sum_{j\ge1}\|L_j-J\|\le
 20\sum_{b=1}^{30}2\cdot80\tau\,2b\,6^{21}
 +40\,\frac{3^{-31}}{1-1/3}
 +\frac{2/3}{1-1/3}\frac{2^{-20}}{1-1/2}<10^{-5}.                        \tag{7}
\]
Из (4),(7) выполнена гипотеза (3) для $R=100$ в прямом случае. Для обратного случая каждый элемент ряда предельных норм равен нулю, поэтому $R$ произволен. Итак, при единственном сдвиге суммарный коэффициент по $h\ge dm$ приближается к 1 для прямого обхода и к 0 для обратного. Для обратного сравнения оценка сильнее всякой фиксированной экспоненты на полосе $aN\le m\le bN$. Нигде не переставлялись сумма по всем временам и предел $N\to\infty$ для конечномерных унитарных матриц.
\newpage
\sect{6. Удаление повторного счёта и полный вес теста}
Сложение по $m$ прямым и $m$ обратным сдвигам пока учитывает пару столько раз, сколько сравнений прошло порог. Для $m\ge3$ это разные перестановки временных индексов. Обозначим кратность через $v$. Ошибка замены $v$ на $\mathbf1_{v>0}$ по модулю ограничена суммой модулей амплитуд по попарным пересечениям, поскольку $0\le v-\mathbf1_{v>0}\le\binom v2$.

Положим $k=\lceil dm\rceil$. Фиксируем два сравнения и по $k$ удовлетворённых уравнений в каждом. Получаем двудольный граф на $2m$ позициях; его рёбра отождествляют буквы двух слов. Полные сравнения --- два совершенных паросочетания. Если их относительная перестановка есть ненулевой поворот, их объединение состоит из циклов одинаковой длины $2m/g\ge4$, $g$ --- НОД величины сдвига и $m$. Если это отражение, имеется не более двух циклов из двух параллельных рёбер (неподвижные индексы), а остальные имеют длину 4. После выбора $2k$ рёбер число целых циклов не превосходит $k/2+1$. Ранг отождествлений равен $2k$ минус число целых циклов. Поэтому свободных букв не более $2m-3k/2+1$.

Наборов индексов не больше $4^m$ на пересечение. По (5) при $j=1$ имеем $\max|U_{XY}|^2\le K/N$ с постоянной $K$. Абсолютная ошибка удаления повторов не больше
\[
             4m^2(4K)^m N^{m-3k/2+1}.                                      \tag{8}
\]
При $d>2/3$ и $aN\le m\le bN$ это меньше $e^{-L m}$ для любого фиксированного $L>0$ начиная с некоторого $N$. Для сравнения только поворотов длина полного цикла известна точно. Для простого $m$ хватает любого фиксированного $d>1/2$, так как целых циклов не больше одного. Равномерно по всем $m$ оценку числа свободных букв нельзя улучшить одним таким подсчётом: при чётном $m$ и относительном сдвиге $m/2$ можно выбирать целые четырёхциклы, сохраняя равенство в ранговой оценке при чётном $k$.

Сумма по сдвигам по (3) равна $m+O(mN^2 e^{-c_0m})$. Вместе с (8) это даёт (1). Циклический вклад равен $D_{m,N}=\sum_x d(x)|a(x)|^2$, где $d(x)$ --- минимальный период. Отличие от $m\Tr[(|U|^{\circ2})^m]$ (здесь берётся след матричной степени матрицы вероятностей) не превосходит $m^2K^mN^{-m/2}$ по подсчёту слов собственного периода. Из (4),(7) при $j=1$ имеем $|\Tr[(|U|^{\circ2})^m]-1|\le N(1/2)^m$ для больших $N$. Это оценивает $D_{m,N}$ с требуемой экспоненциальной точностью.

Обозначим через $H_{m,N}$ однократно учтённую сумму в (1). Для полного веса исходной задачи получено
\[
 \frac1{N^2}\sum_{m=1}^{2N-1} w(m/N)H_{m,N}\longrightarrow
                         \int_0^2 t w(t)\,dt.                              \tag{9}
\]
Действительно, после деления на $N^2$ ошибки (3) с суммой сдвигов мажорируются суммируемой экспонентой с множителем $O(m)$, а ошибки (8) суммируются в $o(1)$ при $d>2/3$. При фиксированном $m$ каждый коэффициент с хотя бы одним совпадением есть сумма $O_m(1)$ следов произведений сжатий с рангом не больше $N$, следовательно имеет порядок $O_m(N)$. Случаи $m=1,2$ учитываем напрямую, когда обратные перестановки совпадают с прямыми. Теорема о доминированной сходимости и сумма Римана дают (9).
\newpage
\sect{6а. Совпадение вне одного окна}
Пусть $m/N\to t>0$, $1\le L<m/2$. Сначала фиксируем начало окна у $x_0$ и прямое совмещение $y$ с $x$. Разрешаем заменить цепочку из $L$ рёбер $x_0,\ldots,x_L$ другой цепочкой с теми же концами. Все остальные рёбра совпадают на тех же местах. Обозначим сумму амплитуд через $f_L$, её часть с одинаковым мультимножеством направленных рёбер через $q_L$, остальную через $f_L-q_L$. Поскольку вне окна рёбра совпадают, условие на мультимножество можно проверять по полным словам.

Пишем $B_{xy}=|U_{xy}|^2$, $J_{xy}=1/N$, $Z=\Tr B^m$. Ниже используются только унитарность $U$ и оценки исходной постановки (7.8): $\max B_{xy}\le K/N$, $\|B^j-J\|\le2^{-j}$. В частности $B^j_{xy}\le K/N$ при $j\ge1$. Сумма по обоим внутренним путям и общему внешнему пути точна:
\[
 f_L=\Tr[B^{m-L} C_L],\qquad (C_L)_{xy}=|(U^L)_{xy}|^2,\qquad
 |f_L-1|\le N2^{-(m-L)}.                                                    \tag{12}
\]
Дважды стохастическая матрица $C_L$ сжимает евклидову норму. Слагаемые $q_L$ положительны или равны нулю, поскольку при перестановке направленных рёбер амплитуда сохраняется. Введём $S(u)=\sinh(u)/u$ с $S(0)=1$ и $F(b)=S(b^2/2)$.

\textbf{Теорема.} При $L/\sqrt N\to b\ge0$ имеем
\[
                q_L\longrightarrow F(b),\qquad f_L-q_L\longrightarrow1-F(b).
                                                                            \tag{13}
\]
Пусть теперь разрешено любое прямое циклическое совмещение $y$, а начало окна у $x$ по-прежнему фиксировано. Обе суммы по множествам пар, делённые на $m$, имеют те же пределы, что в (13). Действительно, при двух допустимых сдвигах совпадает не менее $m-L$ позиций в каждом сравнении. При $L=o(m)$ ошибка перехода к объединению по (8) меньше $e^{-A m}$ для любого фиксированного $A>0$. Это верно и для каждой из двух частей по отдельности. При $L=\lfloor\sqrt m\rfloor$ в (13) следует подставлять $b=\sqrt t$.

\subt{Путь с отмеченным ребром} Положим $Q(x)=\prod_i B_{x_{i+1},x_i}$. К отрезку добавим отличимое искусственное ребро $x_L\to x_0$. Пронумеруем $L+1$ рёбер по обходу; $\sigma$ переводит каждую метку в следующую. На множестве меток одного начала разрешена произвольная перестановка выходящих рёбер. Для совокупной перестановки $\pi$ следующий переход имеет метку $\pi\sigma(i)$. Один обход эйлерова графа получается тогда и только тогда, когда $\pi\sigma$ есть один цикл. Его начало однозначно отмечено искусственным ребром. Если $T(x)$ есть число таких перестановок, а $n_{ab}$ --- кратности реальных рёбер, то
\[
                       q_L=\sum_x Q(x)\frac{T(x)}{\prod_{ab} n_{ab}!}.       \tag{14}
\]
Действительно, каждое результирующее слово допускает ровно $\prod n_{ab}!$ маркировок реальных вхождений. Искусственное ребро имеет собственную метку даже при совпадении концов с одним из реальных рёбер.
\newpage
\sect{6б. Подсчёт перестановок и предел окна}
Маргинальный вес отрезка под $Q$ получается из стационарной цепи Маркова с равномерным началом умножением на $N(B^{m-L})_{x_0,x_L}=1+O(N2^{-(m-L)})$. Этот множитель также ограничен $K$ при $m-L\ge1$. Для ограничений равенства на позициях $0,\ldots,L$ обозначим через $h$ число вершин минус число компонент графа равенств. Сумма весов события под $Q$ не превосходит $K(K/N)^h$. При последовательном раскрытии цепи каждая позиция, кроме первой в своём классе, должна принять уже указанное значение.

Пусть $n=L+1$. При перечислении непересекающихся циклов перестановки суммарный вес мажорируется величиной
\[
 K\exp\!\left(\sum_{j=2}^{n} n^j(K/N)^{j-1}/j\right).
\]
Для $L=O(\sqrt N)$ сумма вкладов перестановок хотя бы с одним циклом длины $\ge3$ равна $O(N^{-1/2})$ по той же экспоненциальной мажоранте. Хвост по числу транспозиций равномерно мал по хвосту экспоненциального ряда. Для фиксированного числа $r$ транспозиций второй момент числа допустимых наборов ограничен: граф двух наборов на $v$ различных индексах имеет $c\le v/2$ компонент и общий вклад $O_r(n^v N^{-(v-c)})=O_r(1)$. Вероятность повторения реального ребра есть $O(L^2/N^2)$, так как любая пара равных направленных рёбер имеет ранг два, в том числе в соседних позициях. Неравенство Коши для ограниченного $r$ и затем оценка хвоста позволяют заменить знаменатель в (14) единицей с ошибкой $o(1)$.

При $b>0$ и данном $r$ все $2r$ концов различны. Число выборов с промежутком между соседними концами меньше $12\log N$ есть $O_r(n^{2r-1}\log N)$. Их общий вклад $o(1)$. В прочих выборах произведение $2r$ переходных вероятностей на цикле равномерно равно $N^{-2r}(1+o(1))$, поскольку и внешний промежуток длинный. Сумма по $r$ свободным меткам вершин даёт $N^{-r}(1+o(1))$. При $b=0$ экспоненциальная мажоранта уже даёт исчезновение нетождественных вкладов.

После стягивания неизменённых участков остаётся равномерное спаривание $2r$ концов на круге. Доля спариваний, сохраняющих один обход, равна $p_r=0$ для нечётных $r$ и $1/(r+1)$ для чётных $r$. Действительно, на стандартном представлении группы $S_{2r}$ индикатор одного цикла равен $\det(I-\sigma\pi)/(2r)$. У одного цикла нетривиальные корни единицы дают определитель $2r$; у остальных есть неподвижный вектор. Внешние степени стандартного представления неприводимы [4]. Усреднение по классу спариваний и лемма Шура дают произведение характеров, делённое на размерность $\binom{2r-1}k$. У длинного цикла $\chi_k(\sigma)=(-1)^k$, а $\sum_k\chi_k(\pi)z^k=(1-z^2)^r/(1+z)$. Заменяя $[2r\binom{2r-1}k]^{-1}$ на $\int_0^1 x^k(1-x)^{2r-1-k}dx$, получаем $p_r=\int_0^1(1-2x)^r dx$. Нулевой случай имеет вес 1. Число спариваний с выбором концов равно $\binom n{2r}(2r-1)!!$; после умножения на $N^{-r}$ оно стремится к $\mu^r/r!$, $\mu=b^2/2$. Сумма $\sum_{r\ge0}p_r\mu^r/r!=S(\mu)$ доказывает (13).
\newpage
\sect{6в. Свободное положение окна и кубический корень}
Теперь разрешим одновременно \emph{любое} начало окна в $x$ и любой прямой сдвиг $y$. Включаем каждую подходящую пару один раз. Обозначим её класс через $\mathcal W_L$, суммы для одинаковых и различных мультимножеств через $A_L$ и $H_L$ соответственно. Пусть $L=\lfloor\lambda m^{1/3}\rfloor$ при фиксированном $\lambda>0$, $m/N\to t>0$. Тогда
\[
 \boxed{\displaystyle
     \frac{A_L}{m}\to 1+\frac{\lambda^3 t^2}{6},\qquad
     \frac{H_L}{m}\to-\frac{\lambda^3 t^2}{6}.}                            \tag{15}
\]
В частности, при $L=\lfloor m^{1/3}\rfloor$ и $m/N\to1$ отрицательная сумма равна $-m/6+o(m)$. Закон справедлив по всем нечётным $N$, включая случаи, когда $m$ составно.

\subt{Точное устранение метки начала} При одном совмещении несовпавшие внутренние вершины должны лежать в $L-1$ последовательных позициях. Непустой набор, помещающийся в такой интервал при $L<m/2$, имеет единственный кратчайший содержащий его интервал: внешний циклический промежуток строго длиннее всех внутренних. Если минимальное допустимое число рёбер $\ell$, то число возможных начал есть $L-\ell+1$. Тождественная пара имеет $m$ возможных начал при любом $L\ge1$. При вычитании подсчётов для $L$ и $L-1$ каждая уже допустимая нетождественная пара остаётся с коэффициентом один. Инвариантность циклической суммы даёт при одном совмещении точные ответы
\[
    Z+m(q_L-q_{L-1}),\qquad Z+m(f_L-f_{L-1})                                \tag{16}
\]
для одинаковых мультимножеств и для всего класса соответственно ($2\le L<m/2$). В частности, $q_L-q_{L-1}$ не убывает по $L$, поскольку положительные классы вкладываются друг в друга. Это даёт выпуклость, необходимую для вычисления разностей. При переходе к объединению по сдвигам умножаем (16) на $m$ с ошибкой (8); число положений окна в этой оценке не появляется.

\subt{Первое отличное от нуля слагаемое} При $\log N\ll L\ll\sqrt N$ и $m/N\to t>0$ равномерно в смысле последовательностей
\[
                          q_L-Z\sim L^4/(24N^2).                            \tag{17}
\]
Одна транспозиция расщепляет обход. Из трёх паросочетаний на четырёх отмеченных позициях ровно перекрёстное сохраняет его. Их суммарный вес равен $\binom{L+1}{4}N^{-2}(1+o(1))$: короткие промежутки вносят $O(L^3\log N/N^2)$. По мажоранте раздела 6б циклы длины хотя бы три дают $O(L^3/N^2)$, а три и более транспозиции дают $O(L^6/N^3)$. Это верхняя оценка (17). Для нижней оставляем четыре хорошо разделённые позиции, требуем ровно два назначенных совпадения и запрещаем все дополнительные равенства. Совместно с дополнительным равенством ранг не менее трёх; если оно затрагивает ноль, один или два новых индекса, число выборов $O(L^4+L^5+L^6)$. Вклад исключённых случаев $o(L^4/N^2)$. В оставшихся словах повтор направленного ребра невозможен, поскольку концы двух пар разнесены более чем на один шаг. Перекрёстная замена даёт второе различимое слово. Это доказывает нижнюю оценку с тем же коэффициентом.
\newpage
\sect{6г. Скорость роста и исходный вес}
По выпуклости разность $q_L-q_{L-1}$ заключена между наклонами хорд до $\lfloor(1-\epsilon)L\rfloor$ и $\lceil(1+\epsilon)L\rceil$. Применяя (17) при всех трёх аргументах и затем устремляя $\epsilon\downarrow0$, получаем $q_L-q_{L-1}\sim L^3/(6N^2)$. Подстановка в (16) и оценка (12) доказывают (15). Тот же аргумент из (13) при $L/\sqrt N\to b>0$ даёт второй закон для окон с произвольным началом:
\[
   \frac{\sqrt N}{m^2} A_L\to F'(b),\qquad
   \frac{\sqrt N}{m^2} H_L\to -F'(b),\qquad
   F'(b)=b\frac{\mu\cosh\mu-\sinh\mu}{\mu^2},\quad \mu=b^2/2.              \tag{18}
\]
Таким образом, при $m\asymp N$ каждая из двух сумм достигает порядка $N^{3/2}$, тогда как сумма обеих равна $m+o(1)$. Начало окна и совмещение слов нигде в (15),(18) не являются метками пары.

\subt{Взвешивание} Для определённости на $m<100$ полагаем рассматриваемые оконные классы пустыми. Для кубического закона ниже берём $\lambda=1$. Из (13) и (15) для части с разными рёбрами, в нормировке $N^{-2}\sum_{m<2N}w(m/N)$, получаем соответственно
\[
 -J_*=-\int_0^2 t w(t)[S(t/2)-1]dt,\qquad
                   -\frac16\int_0^2t^3 w(t)\,dt .                           \tag{19}
\]
В первом интеграле начало фиксировано, $L=\lfloor\sqrt m\rfloor$; во втором оно свободно, $L=\lfloor m^{1/3}\rfloor$. Для обоснования у нуля при фиксированном начале используем $|f_L|\le K$ и равномерную экспоненциальную мажоранту $q_L=O(1)$. При свободном начале числитель положительной прибавки к $Z$ в (16) ограничивается по выпуклости через $m(q_{2L}-Z)/L=O(m(L+1)^3/N^2)$ из верхней комбинаторной оценки (17) без отсечения коротких промежутков. В полном классе $|f_L-f_{L-1}|\le 2N2^{-(m-L)}$; для $m\le C\log N$ вместо этого достаточно $|f_L|\le K$, что после суммирования $O((\log N)^3/N^2)$. При больших $m$ применяем экспоненциальную оценку, затем (8) с $d=0.9$; суммарная нормированная ошибка стремится к нулю. На $m/N\ge\epsilon>0$ достаточно равномерной сходимости, а область $m/N<\epsilon$ даёт $O(\epsilon^2)+o(1)$.

Для исходного $w=(\phi*\phi)^2/(4I^4)$ и $\delta_*=\int_1^2(t-1)w(t)dt$ точная рациональная оценка из архива даёт
\[
                       2.99343<J_*/\delta_*<2.99854,\quad 11.87764<\frac{\int_0^2 t^3 w(t)\,dt}{6\delta_*}<11.89788.
\]
В программе интервалы для $\phi$ получены монотонностью на каждой половине носителя и рядами Тейлора для $\exp(-y/64)$ с шестью возведениями в квадрат. На сеточных узлах свёртка ступенчатых границ вычисляется целочисленным умножением полиномов. На каждой следующей ячейке используются два соседних узла: свёртка симметричных унимодальных функций не возрастает на положительной полуоси (разложение по индикаторам симметричных интервалов). Общая нормировка сократилась. Противоположная положительная прибавка в каждом из (19) равна модулю выписанного отрицательного интеграла. Для вывода (P) остаётся оценить пары вне оконного класса.
\newpage
\sect{6д. Вероятность существования второго обхода}
У среднего (13) есть распределительная интерпретация. Выберем $x$ с вероятностью $Q(x)/Z$ и фиксируем окно $L/\sqrt N\to b>0$. Вероятность того, что разрешена хотя бы одна отличная от исходной внутренняя цепочка с теми же рёбрами, стремится к
\[
                 1-e^{-\mu}\sum_{r=0}^\infty
                         \frac{(2\mu)^r}{r!(r+1)!},\qquad \mu=b^2/2.        \tag{20}
\]
Случайность здесь означает лишь выбор исходного пути с положительным весом, матрица $U_N$ фиксирована.

Действительно, отметим пару $i<j$ при $x_i=x_j$ точкой $(i/n,j/n)$ в треугольнике $0<u<v<1$, $n=L+1$; граничные точки можно отбросить с вероятностью $1-o(1)$. Получаем в пределе пуассоновскую меру интенсивности $b^2\,du\,dv$. Проверим по смешанным факториальным моментам чисел точек в конечном наборе непересекающихся прямоугольных областей с границей меры нуль, обрезанных треугольником. Для $k$ различных пар с $v'$ различными концами и $c$ компонентами ранг равен $v'-c$. Из раздела 6б вклад каждого типа графа есть $O_k(n^{v'}N^{-(v'-c)})$. При общей вершине $c<v'/2$, поэтому он исчезает. Случаи с близкими, но разными индексами исчезают после отсечения $12\log N$. На остальных наборах равномерный вес $N^{-k}(1+o(1))$ оставляет произведение обычных римановых интегралов. Первый момент обеспечивает тесноту распределений счётчиков, а независимые пуассоновские пределы определяются моментами. Формула для всего треугольника даёт полную массу. Вероятность тройных совпадений или повторных рёбер стремится к нулю по оценкам ранга. У предельной конечной конфигурации почти наверное все концы различны и не лежат на границе. Поэтому её порядок концов переносится в предел (разбиение мелкой сеткой и затем снятие сетки).

Условно по числу $r$ точек имеем равномерное спаривание $2r$ позиций. Если есть две перекрёстные хорды, соответствующие две транспозиции дают второй обход. Если все хорды неперекрёстны, любое непустое подспаривание имеет хорду с соседними концами на сжатом круге. После соответствующей транспозиции одна из двух позиций становится неподвижной для $\pi\sigma$. Единого цикла нет. Число неперекрёстных спариваний равно $C_r=\binom{2r}r/(r+1)$: партнёр первого конца разбивает круг на две независимые чётные части, откуда рекурсия Каталана. Итак, условная вероятность единственности равна $C_r/(2r-1)!!=2^r/(r+1)!$ (при $r=0$ равна 1). Усреднение по пуассоновскому $r$ даёт (20).

\subt{Выше масштаба $\sqrt N$} При $\sqrt N\ll L=o(m)$ уже для одного начала и совмещения
\[
                          \frac{\sqrt N}{L}\log q_L\to+\infty.              \tag{21}
\]
Разместим $r=\lfloor L/(\lfloor b\sqrt N\rfloor+g)\rfloor$ блоков длиной $\lfloor b\sqrt N\rfloor$, разделённых $g=\lceil12\log N\rceil$ обычными шагами. Независимый выбор перестановки рёбер в каждом блоке даёт различные глобальные слова, если различаются результирующие внутренние слова. Пусть $K_b(x,y)$ --- сумма весов путей блока от $y$ к $x$, умноженных на число партнёров каждого пути. Среднее $N^{-1}\sum_{xy}K_b(x,y)\to F(b)$ по (13) и равномерной оценке $N(B^{m-\lfloor b\sqrt N\rfloor})_{xy}=1+o(1)$. Все связующие матрицы, включая внешнюю, поэлементно не меньше $(1-N2^{-g})J$. След произведения не меньше $[(1-N2^{-g})N^{-1}\sum K_b]^r$. Нижний предел в (21) тем самым $\ge b^{-1}\log F(b)$ при любом фиксированном $b>0$. Правая часть неограниченно растёт при $b\to\infty$.
\newpage
\sect{6е. Общая непрерывная дуга длиннее половины обхода}
Для непрерывной общей дуги при прямой ориентации можно улучшить порог доли совпадений при суммировании по всем сдвигам. Пусть $0<\eta<1/2$, $aN\le m\le bN$, $2\le L\le (1/2-\eta)m$. Тогда при подсчёте каждой пары один раз
\[
                     \sum_{(x,y)\in\mathcal W_L} a(x)\overline{a(y)}
                                  =m+O_{a,b,\eta}(e^{-c m}).                \tag{22}
\]
Это охватывает изменения внутри блока, длина которого пропорциональна полному времени $m$. Доказательство опирается только на одношаговые оценки (7.8). Для фиксированного совмещения (12),(16) дают главный член $Z$ с ошибкой $O(mN2^{-(m-L)})$. Остаётся оценить пересечения разных совмещений, когда все положения окон уже разрешены.

Пусть $k=m-L$. Выберем по одному отрезку из $k$ последовательных совпавших позиций при двух сдвигах. В графе отождествлений раздела 6 наложение полных сравнений состоит из $g=\gcd(m,s_1-s_2)\le m/2$ циклов. Их проекции на первую копию слова --- классы вычетов по модулю $g$. Среди $m-k$ последовательных позиций вне первого отрезка встречаются ровно $\min(g,m-k)$ таких классов. Совпадения на отрезках поэтому содержат не более $\max(0,g-m+k)\le k-m/2$ полных циклов. Ранг $2k$ выбранных равенств не меньше $2k-(k-m/2)=k+m/2$. При фиксированных отрезках и сдвигах остаётся не более $2m-(k+m/2)$ свободных букв. После умножения на мажоранту амплитуды $(K/N)^m$ и суммирования не более чем по $m^4$ выборам получаем абсолютную ошибку
\[
                            m^4 K^m N^{m/2-k}.                              \tag{23}
\]
Она меньше любой заранее выбранной экспоненты $e^{-Am}$ при $k\ge (1/2+\eta)m$ и больших $N$. Кратность $h$ переучтённой пары исправляется оценкой $h-1\le\binom h2$. Это доказывает (22). Для $L=1$ класс состоит из циклических копий и вывод следует из оценки на $D_{m,N}$.

При $L/m\le 1/2-\eta$ и $L/\sqrt N\to\infty$ доказательство (21) остаётся верным: внешняя дуга по-прежнему длиннее $g$. Поэтому суммарный положительный вклад одинаковых рёбер уже внутри этого класса больше $m\exp(C L/\sqrt N)$ при любой фиксированной $C$ и достаточно больших $N$. Остальные пары внутри него компенсируют эту величину до $m+O(e^{-cm})$ по (22). Для исходного парного теста требуется вычислить вклад пар, не содержащих такой общей дуги.
\newpage
\sect{6ж. Закон кратчайшего участка перестановки рёбер}
Пусть $x$ выбрано с вероятностью $Q(x)/Z$. Совмещение фиксировано. Определим $D(x)$ как минимальное число последовательных рёбер циклического слова, которые можно переставить в отличный путь с теми же концами, не меняя остальную часть слова. При невозможности полагаем $D=+\infty$. Для каждого $u>0$ и $m/N\to t>0$
\[
  \boxed{\displaystyle
  \Pr\!\left\{D(x)>u(N^2/m)^{1/3}\right\}\longrightarrow e^{-u^3/6}.}       \tag{24}
\]
Более точно, при $L=\lfloor u(N^2/m)^{1/3}\rfloor$ число различных нетождественных слов-партнёров, получаемых заменой внутри одного произвольно расположенного окна длины $L$, сходится по распределению к $\operatorname{Pois}(u^3/6)$. В обоих утверждениях исходная квантовая карта детерминирована; вероятность относится к положительному весу отдельно взятого исходного пути. Это утверждение о траекториях, для статистики собственных фаз нужен дальнейший переход.

\subt{Перекрёстные пары в короткой дуге} Считаем наборы четырёх позиций $a<b<c<d$ в циклическом порядке, при которых $d-a\le L$, $x_a=x_c$, $x_b=x_d$. Первой считаем позицию после большого внешнего промежутка. При $L<m/2$ такой выбор единственен; индексы временно поднимаются в целые. Всего возможны $m\binom L3$ наборов. Обозначим число реализованных через $X_L$. У хорошо разделённых четырёх концов (все промежутки $\ge12\log N$) вероятность равна $N^{-2}(1+o(1))$. Для любого набора она $O(N^{-2})$: циклическая маргиналь на $v$ позициях есть произведение $v$ переходных вероятностей, делённое на $Z$. При $c'$ независимых метках суммирование даёт $O_v(N^{-(v-c')})$. Число наборов с коротким промежутком есть $O(m L^2\log N)$. Поэтому $\mathbb E X_L\to u^3/6$.

Для факториального момента порядка $h$ выберем $h$ различных четвёрок. Объединим в группы четвёрки, чьи начальные позиции находятся друг от друга на круговом расстоянии $\le2L$, с транзитивным замыканием. Для $z$ групп и $v$ различных позиций имеется $O_h(m^z L^{v-z})$ конфигураций данного типа, поскольку каждая группа расположена на дуге длиной $O_h(L)$. Граф равенств на $v$ позициях имеет $c'\le v/2$ компонент. Суммарный вклад типа не превосходит
\[
 O_h(m^z L^{v-z}N^{-(v-c')})=O_h(N^{2z/3-v/6}).
\]
Каждая группа содержит не менее четырёх позиций. При двух разных четвёрках в одной группе их объединение содержит как минимум пять: четыре заданные позиции допускают только одно перекрёстное спаривание. Тогда $v\ge4z+1$ и вклад исчезает. Когда $z=h$, группы разнесены, позиции в разных группах различны. Удаление наборов с коротким внутренним промежутком даёт $o(1)$ по той же одночетвёрочной оценке, умноженной на ограниченный вклад остальных наборов. Число остальных упорядоченных конфигураций асимптотически равно $[m\binom L3]^h$ (сближение двух начал отбрасывает долю $O_h(L/m)$). Их общий вес на конфигурацию $N^{-2h}(1+o(1))$ по смешиванию. Получили все факториальные моменты пуассоновского закона. Он определяется моментами, откуда $X_L\Rightarrow\operatorname{Pois}(u^3/6)$.
\newpage
\sect{6з. От совпавших позиций к различным траекториям}
Назовём четвёрку из предыдущего раздела изолированной, если в её дуге $[a,d]$ нет никаких равенств, кроме двух назначенных, промежутки её концов $\ge12\log N$ и начало никакой другой реализованной четвёрки не находится на расстоянии $\le2L$ от её начала. Пусть $X_L^0$ есть число таких четвёрок. Тогда
\[
                       \mathbb E(X_L-X_L^0)=o(1).                           \tag{25}
\]
Действительно, короткие промежутки уже оценены, а ожидаемое число упорядоченных пар близких четвёрок стремится к нулю по доказательству второго факториального момента. Для дополнительного равенства в дуге ранг трёх условий равен трём: оно не является повтором одной из двух назначенных пар. Общее число различных индексов лежит между четырьмя и шестью. При $v$ индексах число конфигураций не больше $O(m L^{v-1})$, так что вклад $O(m(L^3+L^4+L^5)/N^3)=o(1)$. Это доказывает (25).

У изолированной четвёрки имеются три последовательных участка $P,Q_1,R$ с концами $A\to B$, $B\to A$, $A\to B$. Заменим $P Q_1 R$ на $R Q_1 P$. Обе последовательности имеют ту же длину и те же концы, и первая внутренняя вершина у них различна: иначе в позициях $a+1$ и $c+1$ было бы дополнительное равенство. Остальной путь сохранён. Дуги изолированных четвёрок попарно не пересекаются, поэтому выполненные по одной замены дают различные слова. Если $T_L$ есть фактическое число нетождественных партнёров при одном совмещении и любом начале окна, получаем поточечное неравенство $T_L\ge X_L^0$.

С другой стороны, положительная формула (16) даёт точное ожидание
\[
                  \mathbb E T_L=\frac{m(q_L-q_{L-1})}{Z}\to u^3/6.
\]
Здесь использованы (17), выпуклость и $Z\to1$. Но и $\mathbb E X_L^0\to u^3/6$ по (25). Следовательно, $\mathbb E|T_L-X_L^0|\to0$. Целочисленность и неравенство Маркова дают $\Pr(T_L=X_L)\to1$; получаем заявленный пуассоновский предел. Поскольку $D>L$ равносильно $T_L=0$, следует (24). Циклические копии исходного слова, отличные при данном совмещении, также можно удалить: такая пара наряду с данным совмещением допускает другое, полное; её суммарный положительный вес оценивается (8).

\subt{Конечная проверка подсчёта} Для независимых равномерных букв при фиксированном $n=L+1\ge3$ коэффициент при $N^{-2}$ в $q_L-1$ равен $\binom{n+1}{4}-\binom{n-1}{2}$. Одна тройная вершина допускает перестановку двух экскурсий, кроме случая трёх соседних позиций (оба ребра --- одинаковые петли). При двух двойных вершинах подходит перекрёстная пара, кроме случая двух одинаковых однорёберных переходов $A\to B$. Получаем $\binom n3-(n-2)+\binom n4-\binom{n-2}2$. Все события более высокого ранга имеют порядок $O_n(N^{-3})$. Полный перебор типов равенства при $n\le9$ и независимая точная проверка (12),(16) по 5\,045\,113 парам слов приведены в архиве.

При $m/N\to1$ медиана $D$ асимптотически равна $(6\log 2)^{1/3}N^{1/3}$. В масштабе $\sqrt N$ для одного фиксированного начала вероятность существования второй цепочки дана формулой (20); в масштабе $(N^2/m)^{1/3}$ для всех начал её задаёт (24). Число возможных начал и устранение их повторов входят в определение величины. Вычисленный закон описывает, когда две траектории впервые могут дать тождественные произведения матричных элементов посредством локальной перестановки направленных рёбер.
\newpage
\sect{7. Что ограничения означают для собственных состояний}
Пусть $\Pi_\theta$ --- спектральный проектор соответствующего двойного оператора, в прямом случае на $e^\perp$. Для единичного $v\in\mathrm{ran}\,P_0$ и целого $T$ имеем
\[
 \|\Pi_\theta v\|^2\le\left\|T^{-1}\sum_{j=0}^{T-1} e^{-ij\theta}V^jv\right\|^2
 \le T^{-1}+\frac2{T^2}\sum_{j=1}^{T-1}(T-j)\|K_{j,N}\|.
\]
По (4),(7) верхняя грань по всем углам и таким $v$ имеет верхний предел не больше $(1+2\cdot10^{-5})/T$ в прямом случае и $1/T$ в обратном. Увеличиваем $T$. Та же оценка с множителем $1+o(1)$ при фиксированном $T$ применима к проекциям на произвольные дуги длины $o(1)$: на такой дуге модуль нормированного геометрического ряда, центрированного в середине дуги, равномерно близок к единице.

В прямом случае векторы $u_j\otimes\bar u_j$ ортонормированы и принадлежат собственному пространству угла нуль до удаления $e$. Для диагонального $v=\sum v_x|x,x\rangle$, $\sum v_x=0$, суммы квадратов перекрытий с ними оцениваются найденной границей: его проекция на $e$ уже равна нулю. Эти перекрытия суть координаты матрицы $D^*v$, $D_{xj}=|u_j(x)|^2$. Так как $D$ дважды стохастична, получается первая формула (2). В частности для любых, в том числе выбранных по спектру, множеств позиций $E$ и номеров состояний $B$,
\[
 \left|\sum_{x\in E,\,j\in B}|u_j(x)|^2-\frac{|E||B|}{N}\right|
                         \le o(1)\sqrt{|E|\,|B|}.
\]

У обратного оператора $u_j\otimes u_j$ ортонормированы и имеют угол нуль. Теперь диагональный $v$ произволен. Отсюда вторая формула (2). В частности, для любой последовательности фаз $|h_x|=1$, даже выбираемой после нахождения всего собственного базиса, положим $C_h u(x)=h_x\bar u(x)$. Тогда
\[
             \sup_{|h_x|=1}\frac1N\sum_{j=1}^N
                     |\langle u_j,C_h u_j\rangle|^2\longrightarrow0.       \tag{10}
\]
Действительно, берём $v_x=h_x/\sqrt N$. Если после одного диагонального изменения фаз положительная доля состояний была бы близка к вещественным с точностью до индивидуальной скалярной фазы, их слагаемые были бы близки к 1. Формула (10) это исключает. Вывод равномерен по выбору ортонормированных базисов внутри возможных кратных собственных пространств. Здесь исследуется одна общая позиционная калибровка для всех состояний.

\textbf{Независимая граница с одним шагом.} Часть вывода о прямых парах имеет короткое доказательство даже без многошаговой формулы (4). Пусть $b=\||U|^{\circ2}-J\|$, $0<r<1$. Для веса $z$ за \emph{несовпадение}, $|z|=r$, на $e^\perp$ спектральный радиус $(P_0+zQ)V_+$ не больше
\[
            R_b(r)=\tfrac12\big(b(1-r)+\sqrt{b^2(1-r)^2+4r}\big).
\]
В блоках $V_+=\left(\begin{smallmatrix}D&B\\C&H\end{smallmatrix}\right)$ относительно $P_0,Q$ функция $D+wB(I-wH)^{-1}C$ есть сжатие в единичном круге по унитарности. Лемма Шварца--Пика для скалярных матричных элементов даёт оценку $(b+|w|)/(1+b|w|)$. При собственном значении $|\lambda|>r$ подставляем $w=z/\lambda$ в уравнение на первый блок, что даёт указанное квадратное неравенство. Поскольку $R_b'(1)=(1-b)/2$, оценка Коши подавляет номера несовпадений $\le\alpha m$ при $\alpha<(1-b)/2$. Уже исходная однократная оценка (7.7) даёт $b<1/3$ при больших $N$, а (8) удаляет повторы при доле совпадений $d>2/3$. Многошаговая оценка нужна для всех $d>0$ при фиксированном сдвиге и для (2),(10).
\newpage
\sect{8. Точная простота на двух бесконечных семействах}
Есть прямой результат об исходных собственных фазах, независимый от (1)--(10). Для $N=3^k$ и $5^k$, $k\ge1$, спектр $U_N$ прост. Более того, спектр $U_N^m$ прост для каждого $m$, взаимно простого с $(p+1)p^{k-1}$, $p=3,5$. Это включает степени порядка $N$.

Действие $S$ на метке $(a,b)$ соответствует умножению $a+b\sqrt2$ на $v=3+2\sqrt2$. Для обоих $p$ число 2 --- квадратичный невычет. Примитивный элемент кольца $(\Z/p^l)[\sqrt2]$ обратим: его редукция ненулевая в поле $\mathbb F_{p^2}$, а норма обратима по модулю $p^l$. Поэтому длина каждой примитивной орбиты равна порядку $v$. Этот порядок по модулю $p$ равен $p+1$. Например, $v^2=17+12\sqrt2$ и $v^3=99+70\sqrt2$; непосредственная проверка исключает меньшие делители. В $v^{p+1}-1$ обе координаты делятся на $p$, но не обе на $p^2$ ($v^4=577+408\sqrt2$, $v^6=19601+13860\sqrt2$). По биномиальной формуле возведение в степень $p$ увеличивает общую $p$-адическую нормировку разности с единицей ровно на один. Таким образом порядок по модулю $p^l$ равен $(p+1)p^{l-1}$.

По аддитивным порядкам меток количество всех циклов $S\bmod p^k$ равно
\[
       1+\sum_{l=1}^k\frac{p^{2l}-p^{2l-2}}{(p+1)p^{l-1}}=p^k.
\]
Чистое сопряжение переставляет вейлев базис. Значит, размерность коммутанта $M$ равна числу циклов, то есть $N$. Но она равна также сумме квадратов кратностей, при сумме кратностей $N$. Отсюда простота. При указанном $m$ циклы $S^m$ имеют прежние длины и разбиение, так что тот же вывод верен для $M^m$. Родственные вопросы о линейных картах на степенях простых изучаются в [3].

Дискриминант характеристического многочлена $(e^{itA}M)^m$ есть конечный экспоненциальный полином $\sum_\nu c_\nu e^{it\nu}$ с алгебраическими коэффициентами и алгебраическими частотами: используем спектральные проекторы эрмитовой алгебраической матрицы $A$. После группировки одинаковых частот значение при $t=0$ ненулевое, поэтому по теореме Линдемана--Вейерштрасса ненулево и значение при $t=10^{-80}N$. Это доказывает утверждение для исходного нелинейного оператора.

В частности, на этих двух последовательностях нормированная мера разностей $N^{-1}\sum_{ij,l}\delta_{N(\theta_i-\theta_j+l)}$ имеет в нуле атом в точности 1. При каждом таком $N$ среднее $N^{-1}M_0^{-1}\sum_{m=1}^{M_0}|\Tr U^m|^2$ стремится к 1 при $M_0\to\infty$ по сумме геометрических прогрессий. Для (P) одновременно растут $M_0$ и $N$, поэтому указанное усреднение не вычисляет требуемого плато.
\newpage
\sect{9. Точная положительная формула для абелева теста}
Для $r=e^{-2s/N}$ запишем $L_N(s)=N^{-2}\sum_{m\ge1}r^m|p_m|^2$. Для любого унитарного спектра на окружности установим
\[
 L_N(s)-\frac1{e^{2s}-1}=\frac1{N^2}\sum_{k=1}^{N-2}h_k|p_k|^2
       +\frac1{2N^2D}\sum_{i\ne j}P_r(x_{ij})F_N(x_{ij}),                 \tag{11}
\]
где $x_{ij}=\cos(2\pi(\theta_i-\theta_j))$, $P_r(x)=(1-r^2)/(1-2rx+r^2)$,
\[
 F_N(\cos u)=\left(\frac{\sin(Nu/2)}{\sin(u/2)}\right)^2,\qquad
 D=F_N((1+r^2)/(2r))=\sinh^2 s/\sinh^2(s/N).
\]
Значение при $u=0$ берётся по непрерывности. В правой части все слагаемые неотрицательны. Действительно, $H=P_r(1-F_N/D)$ есть многочлен степени $N-2$: числитель исчезает в полюсе $P_r$. Из $P_r(\cos u)=\sum_{k\in\Z}r^{|k|}e^{iku}$ и коэффициентов $N-|j|$ у $F_N$ получаем для $0\le k\le N-2$
\[
       h_k=D^{-1}\sum_{j=k+1}^{N-1}(N-j)(r^{k-j}-r^{j-k})>0 .
\]
Суммируем $H$ и остаток $P_rF_N/D$ по всем парам. $h_0$ и члены $i=j$ независимы от спектра. Их удобно вычесть с помощью регулярного $N$-угольника, у которого исчезают все $p_k$, $1\le k\le N-2$, а $F_N=0$ на каждой нетождественной разности. Его $L_N=r^N/(1-r^N)$. Это доказывает (11), в том числе при кратных фазах. Равенство возможно только для полного регулярного многоугольника, поскольку все отношения различных по индексу точек должны быть нетождественными корнями степени $N$.

Для данного $U_N$ минимум строг при каждом нечётном $N\ge3$: $\Tr U_N$ ненулевой, так как соответствующий экспоненциальный полином ненулевой уже при $t=0$ ($|\Tr M|=1$). Из оценки первого следа в предыдущей части и предела коэффициента $h_1$ получается явная конечномасштабная поправка
\[
 \liminf_{N\to\infty} N^2\left(L_N(s)-\frac1{e^{2s}-1}\right)
 \ge (1+300\tau)\frac{\sinh(2s)-2s}{2\sinh^2s}.
\]
Предел $h_1$ вычисляется суммой Римана для $2s^2\int_0^1(1-t)\sinh(2st)dt/\sinh^2s$. Масштаб этой нижней границы ещё допускает разные предельные значения $L_N(s)$.

При простом спектре разложение при $s\downarrow0$, $N$ фиксировано, даёт
\[
 L_N(s)=\frac1{2s}-\frac12+
 s\left[\frac1{6N^2}+\frac1{2N^3}\sum_{i\ne j}\csc^2(\pi(\theta_i-\theta_j))\right]
 + O_N(s^3).
\]
Достаточно подставить $P_r=\sinh(2s/N)/(\cosh(2s/N)-x)$, отдельно выделив диагональ. На размерностях из раздела 8 это доказывает первые два члена при исходном кике. Остаток не заявлен равномерным по $N$.
\newpage
\sect{10. Близость к минимуму и оставшаяся интерференция}
У (11) есть критерий непосредственно для слабого закона. Если на какой-либо бесконечной подпоследовательности для хотя бы одного фиксированного $s>0$ выполняется $L_N(s)-1/(e^{2s}-1)\to0$, то локальный процесс на ней сходится к единичной решётке с равномерным сдвигом.

Пусть $d_i$ --- циклические соседние щели в единицах среднего расстояния. В сумме справа (11) оставим соседние пары. При $d_i=t$ на любом компакте по $t$ вес одного члена после умножения на $N$ равномерно стремится к
\[
       \frac{s^3}{2\sinh^2s}\,
       \frac{(\sin(\pi t)/(\pi t))^2}{s^2+\pi^2t^2}.
\]
На $[0,1-\epsilon]$ минимум положителен. Поэтому доля коротких щелей $d_i\le1-\epsilon$ стремится к нулю. Из $N^{-1}\sum_i d_i=1$ следует $N^{-1}\sum_i|d_i-1|=2N^{-1}\sum_i(1-d_i)_+\to0$. При равномерном выборе начала на окружности вероятность попасть в щель $i$ равна $d_i/N$. Это распределение сближается с равномерным по индексам в полной вариации. Внутри выбранной щели положение равномерно. Для равномерного индекса каждая фиксированная группа соседних щелей сходится по вероятности к последовательности единиц по неравенству Маркова и объединению событий. Любое фиксированное компактное окно покрывается конечной группой таких щелей с вероятностью, стремящейся к единице. Получен заявленный слабый предел. Для карты этой работы приближение к минимуму на бесконечном ряду не доказано.

\textbf{Что именно осталось для (P).} Пусть $G_{m,N}=|p_{m,N}|^2-H_{m,N}$ --- вклад всех остальных пар из раздела 6. Из (9) и определения $\delta_*=\int_1^2(t-1)w(t)dt>0$ следует
\[
 \Delta_{2,N}(b_*)=\delta_*+
              N^{-2}\sum_{m=1}^{2N-1}w(m/N)G_{m,N}+o(1).
\]
Искомая поправка $-\delta_*$ целиком должна приходить из пар, не набирающих заданной доли $d>2/3$ совпадений ни в одном из двух направлений. В этом классе всё ещё находятся пары, посещающие те же рёбра в существенно ином порядке. Их совместная знаковая асимптотика на масштабе $N^2$ неизвестна. Ни (10), ни устранение точных кратностей по разделу 8, ни минимум (11) не определяют её.

\textbf{Источники.} [1] M. Zworski, \emph{Semiclassical Analysis}, гл. 3 и 10, \href{https://doi.org/10.1090/gsm/138}{Graduate Studies in Mathematics 138}. [2] J. H. Hannay, M. V. Berry, \emph{Quantization of linear maps on a torus}, \href{https://doi.org/10.1016/0167-2789(80)90026-3}{Physica D 1 (1980), 267--290}. Для нечётного сектора в разделе 4 дано явное тождество четырёх ветвей. [3] D. Kelmer, \emph{On matrix elements for the quantized cat map modulo prime powers}, \href{https://arxiv.org/abs/0802.3237}{arXiv:0802.3237}. [4] W. Fulton, J. Harris, \emph{Representation Theory. A First Course}, предложение 3.12, \href{https://doi.org/10.1007/978-1-4612-0979-9}{GTM 129}. Обозначения BGS-F и (P) взяты из постановки в архиве. Оценки исходной карты из прежних частей указаны в тех местах, где используются.

\textbf{Проверки.} В архиве находятся целочисленные проверки фаз, орбит и подсчёта ранга, точная рациональная проверка (11), а также независимый перебор пар слов с матричной проверкой коэффициентов для обеих ориентаций. Численный перебор относится к малым размерностям; предельные оценки следуют из приведённых аналитических доказательств.
\end{document}
